\documentclass[tikz,border=4pt]{standalone}
\usepackage{amsmath}
\usepackage{amssymb}
\usepackage{tikz}
\usetikzlibrary{arrows.meta,positioning,calc}

\begin{document}
\begin{tikzpicture}[
    font=\small,
    box/.style={rounded corners=12pt, very thick, minimum width=3.9cm, minimum height=6.1cm, inner sep=0pt, draw},
    title/.style={font=\bfseries\large, text=white, align=center},
    content/.style={font=\small, align=left, text=black}
]

% Colors
\definecolor{blueTop}{RGB}{48,115,170}
\definecolor{blueLight}{RGB}{232,240,248}
\definecolor{redTop}{RGB}{183,61,54}
\definecolor{redLight}{RGB}{249,231,228}
\definecolor{greenTop}{RGB}{77,139,72}
\definecolor{greenLight}{RGB}{232,244,229}
\definecolor{purpleTop}{RGB}{123,66,142}
\definecolor{purpleLight}{RGB}{244,232,247}

% A robust arrow macro: no shapes.arrows needed
\newcommand{\blockarrow}[3]{%
  \begin{scope}[shift={(#1,#2)}]
    \filldraw[fill=#3!88, draw=#3, very thick]
      (-0.34,-0.23) -- (0.18,-0.23) -- (0.18,-0.42)
      -- (0.63,0) -- (0.18,0.42) -- (0.18,0.23)
      -- (-0.34,0.23) -- cycle;
  \end{scope}
}

% ---------------- Case 1 ----------------
\node[box, draw=blueTop, fill=blueLight] (c1) at (0,0) {};
\fill[blueTop, rounded corners=12pt]
    ([xshift=-1.95cm,yshift=3.05cm]c1.center) rectangle
    ([xshift=1.95cm,yshift=1.75cm]c1.center);
\node[title] at ([yshift=2.45cm]c1.center) {Case 1: Risk Control\\(Calibration)};

\node[content] at ([yshift=0.35cm]c1.center) {
\begin{minipage}{3.35cm}
\centering
$\widehat{R}(\tau)=\dfrac{\sum(\hat{y}_i\neq y_i)}{|S_\tau|}$

\vspace{0.2cm}
$\tau_1^*=\inf\{\tau\in[0,1]\mid \widehat{R}(\tau)\leq \alpha\}$

\vspace{0.18cm}
$(f,g)_{\tau_1^*}(X)=
\begin{cases}
(f(X)>0.5),\\
g(X)\geq \tau_1^*,\\
\text{abstain},\\
\text{otherwise.}
\end{cases}$
\end{minipage}
};

\node[content, anchor=west] at ([xshift=-1.65cm,yshift=-2.35cm]c1.center) {
\begin{minipage}{3.25cm}
\textbf{Objective:} Control empirical risk.\\
\textbf{Limitation:} Fails FNR constraints.
\end{minipage}
};

\blockarrow{2.15}{0}{blueTop}

% ---------------- Case 2 ----------------
\node[box, draw=redTop, fill=redLight] (c2) at (4.45,0) {};
\fill[redTop, rounded corners=12pt]
    ([xshift=-1.95cm,yshift=3.05cm]c2.center) rectangle
    ([xshift=1.95cm,yshift=1.75cm]c2.center);
\node[title] at ([yshift=2.45cm]c2.center) {Case 2: FNR Control\\(Calibration)};

\node[content] at ([yshift=0.35cm]c2.center) {
\begin{minipage}{3.35cm}
\centering
$\mathrm{FNR}(\tau)=\dfrac{\sum(\hat{y}_i=0\wedge y_i=1)}{|\{y_i=1\}|}$

\vspace{0.2cm}
$\tau_2^*=\inf\{\tau\in[0,1]\mid \mathrm{FNR}(\tau)\leq \alpha\}$

\vspace{0.18cm}
$(f,g)_{\tau_2^*}(X)=
\begin{cases}
(f(X)>0.5),\\
g(X)\geq \tau_2^*,\\
\text{abstain},\\
\text{otherwise.}
\end{cases}$
\end{minipage}
};

\node[content, anchor=west] at ([xshift=-1.65cm,yshift=-2.35cm]c2.center) {
\begin{minipage}{3.25cm}
\textbf{Objective:} Control empirical FNR.\\
\textbf{Limitation:} Low coverage, high abstention.
\end{minipage}
};

\blockarrow{6.60}{0}{redTop}

% ---------------- Case 3 ----------------
\node[box, draw=greenTop, fill=greenLight] (c3) at (8.90,0) {};
\fill[greenTop, rounded corners=12pt]
    ([xshift=-1.95cm,yshift=3.05cm]c3.center) rectangle
    ([xshift=1.95cm,yshift=1.55cm]c3.center);
\node[title] at ([yshift=2.32cm]c3.center) {Case 3: Risk \&\\FNR Control\\(Cost-sensitive)};

\node[content, anchor=west] at ([xshift=-1.65cm,yshift=1.00cm]c3.center) {
\begin{minipage}{3.25cm}
\textbf{Approach:} Incorporate Cost-sensitive Learning during training to penalize FNs.
\end{minipage}
};

\node[content] at ([xshift=0.10cm,yshift=-0.10cm]c3.center) {
\begin{tikzpicture}[scale=0.85, font=\small]
\node (obj) at (0,0) {Objective};
\node (risk) at (-1.05,0.85) {Risk};
\node (fdr) at (1.0,0.85) {FDR};
\node (cov) at (1.1,-0.45) {Coverage};
\draw[greenTop, very thick,-{Latex[length=2mm]}] (obj) -- (risk);
\draw[greenTop, very thick,-{Latex[length=2mm]}] (obj) -- (fdr);
\draw[greenTop, very thick,-{Latex[length=2mm]}] (obj) -- (cov);
\draw[greenTop, very thick,-{Latex[length=2mm]}] (risk) -- (cov);
\draw[greenTop, very thick,-{Latex[length=2mm]}] (cov) -- (fdr);
\end{tikzpicture}
};

\node[content, anchor=west] at ([xshift=-1.65cm,yshift=-2.15cm]c3.center) {
\begin{minipage}{3.25cm}
\textbf{Objective:} Balance Risk Control, FDR Control, \& Coverage. Addresses Case 1 FNR failure.
\end{minipage}
};

\blockarrow{11.05}{0}{purpleTop}

% ---------------- Case 4 ----------------
\node[box, draw=purpleTop, fill=purpleLight] (c4) at (13.35,0) {};
\fill[purpleTop, rounded corners=12pt]
    ([xshift=-1.95cm,yshift=3.05cm]c4.center) rectangle
    ([xshift=1.95cm,yshift=1.55cm]c4.center);
\node[title] at ([yshift=2.32cm]c4.center) {Case 4: FNR \&\\Coverage Control\\(Cost-sensitive)};

\node[content, anchor=west] at ([xshift=-1.65cm,yshift=1.00cm]c4.center) {
\begin{minipage}{3.25cm}
\textbf{Approach:} Heavily penalize FNs in Cost-sensitive Learning.
\end{minipage}
};

\node[content] at ([xshift=0.05cm,yshift=-0.05cm]c4.center) {
\begin{tikzpicture}[scale=0.85, font=\small]
\node (heavy) at (-0.9,0.25) {Heavily};
\node (cov) at (1.05,0.85) {Coverage};
\node (fnr) at (1.15,-0.45) {FNR Control};
\draw[purpleTop, very thick,-{Latex[length=2mm]}] (heavy) -- (cov);
\draw[purpleTop, very thick,-{Latex[length=2mm]}] (heavy) -- (fnr);
\end{tikzpicture}
};

\node[content, anchor=west] at ([xshift=-1.65cm,yshift=-2.15cm]c4.center) {
\begin{minipage}{3.25cm}
\textbf{Objective:} Improve Coverage while maintaining FNR Control. Addresses Case 2 low coverage.
\end{minipage}
};

\end{tikzpicture}
\end{document}
